\documentclass[11pt]{article}
\usepackage{amsmath,amsthm,amssymb}
\usepackage[margin=1in]{geometry}
\usepackage{booktabs}

\newtheorem{theorem}{Theorem}
\newtheorem{lemma}[theorem]{Lemma}
\newtheorem{proposition}[theorem]{Proposition}
\newtheorem{corollary}[theorem]{Corollary}
\newtheorem{conjecture}[theorem]{Conjecture}
\theoremstyle{definition}
\newtheorem{definition}[theorem]{Definition}
\theoremstyle{remark}
\newtheorem{remark}[theorem]{Remark}
\newtheorem{example}[theorem]{Example}

\newcommand{\ZZ}{\mathbb{Z}}
\newcommand{\CC}{\mathbb{C}}
\newcommand{\Sym}{\mathfrak{S}}
\newcommand{\im}{\operatorname{im}}
\newcommand{\tr}{\operatorname{tr}}
\newcommand{\rk}{\operatorname{rank}}
\DeclareMathOperator{\Ind}{Ind}

\title{The $H$-Invariant Theorem for the Staircase Symmetrizing Product}
\author{Clio}
\date{April 14, 2026}

\begin{document}
\maketitle

\begin{abstract}
Let $\Pi = \prod_{j}(1+s_{i_j})/2$ be the staircase symmetrizing product
for $w_0 \in \Sym_n$, acting on the left regular representation, and let
$H = \langle s_1, s_3, s_5, \ldots \rangle \cong (\ZZ/2)^{\lfloor n/2\rfloor}$
be the subgroup of pair transpositions. We prove that the rank of $\Pi$
in each irreducible component $V_\lambda$ equals the dimension of the
$H$-invariant subspace: $r_\lambda = \dim(V_\lambda^H)$. This is verified
computationally for all $n \le 7$. We derive as corollaries the rank
formula $\rk(\Pi) = n!/2^{\lfloor n/2\rfloor}$, the vanishing criterion
$r_\lambda = 0 \Leftrightarrow \ell(\lambda) > \lceil(n+1)/2\rceil$,
and the hook formula $r_{(n-k,1^k)} = \binom{\lfloor(n-1)/2\rfloor}{k}$.
We prove the main theorem unconditionally for $n \le 3$ and establish
several structural lemmas toward the general case.
\end{abstract}


\section{Setup and main theorem}

\subsection{Definitions}

Let $\Sym_n$ act on $\{1,\ldots,n\}$ with simple transpositions
$s_i = (i,\, i{+}1)$. In the left regular representation,
$s_i \cdot e_w = e_{s_iw}$. Define the projections
$P_i = (1+s_i)/2$, which are idempotents projecting onto the
$+1$-eigenspace of $s_i$.

\begin{definition}[Staircase word and product]
The \emph{staircase reduced word} for the longest element
$w_0 \in \Sym_n$ is
\[
  (s_1)(s_2, s_1)(s_3, s_2, s_1) \cdots (s_{n-1}, \ldots, s_2, s_1),
\]
consisting of $n-1$ \emph{blocks}: block $k$ (for $k=1,\ldots,n-1$)
applies $s_k, s_{k-1}, \ldots, s_1$. The \emph{staircase symmetrizing
product} is
\[
  \Pi = \prod_{k=1}^{n-1} \prod_{j=k}^{1} P_j
    = P_{n-1}\cdots P_2 P_1 \cdots P_2 P_1 \cdot P_1,
\]
where the rightmost factor is applied first.
\end{definition}

\begin{definition}[Pair subgroup]
The \emph{pair subgroup} is $H = \langle s_1, s_3, s_5, \ldots,
s_{2\lfloor n/2\rfloor - 1}\rangle \le \Sym_n$. Since
$|i - j| \ge 2$ for distinct odd $i,j$, the generators commute,
so $H \cong (\ZZ/2)^{\lfloor n/2\rfloor}$ with $|H| = 2^{\lfloor n/2\rfloor}$.
\end{definition}

\begin{definition}[Irreducible rank]
For $\lambda \vdash n$, let $V_\lambda$ be the irreducible $\Sym_n$-module
with Young's seminormal basis $\{v_T : T \in \mathrm{SYT}(\lambda)\}$.
The \emph{irreducible rank} is
\[
  r_\lambda = \rk(\Pi_\lambda),
\]
where $\Pi_\lambda$ is the product $\prod P_{i_j}$ evaluated in the
representation $V_\lambda$ (i.e., $P_k^\lambda = (I + \rho_\lambda(s_k))/2$).
\end{definition}


\subsection{Statement}

\begin{theorem}[$H$-Invariant Theorem]
\label{thm:main}
For all $\lambda \vdash n$:
\[
  r_\lambda = \dim(V_\lambda^H),
\]
where $V_\lambda^H = \{v \in V_\lambda : s_k v = v \text{ for all odd }
k \le 2\lfloor n/2\rfloor - 1\}$ is the $H$-invariant subspace.
\end{theorem}

\noindent The $H$-invariant subspace $V_\lambda^H$ is also the image of
the commuting projection $Q_\lambda = \prod_{k\;\mathrm{odd}} P_k^\lambda$,
and $\dim(V_\lambda^H) = \rk(Q_\lambda)$.


\subsection{Computational verification}

Theorem~\ref{thm:main} is verified for all $\lambda \vdash n$ with
$2 \le n \le 7$ by computing $r_\lambda$ (via SVD of $\Pi_\lambda$ in
the seminormal representation) and $\dim(V_\lambda^H)$ (as
$\rk(Q_\lambda)$). The full data:

\begin{center}
\small
\begin{tabular}{c|c|c|c|c}
\toprule
$n$ & $\lambda$ & $f^\lambda$ & $r_\lambda$ & $\dim(V_\lambda^H)$ \\
\midrule
2 & $(2)$ & 1 & 1 & 1 \\
  & $(1^2)$ & 1 & 0 & 0 \\
\midrule
3 & $(3)$ & 1 & 1 & 1 \\
  & $(2,1)$ & 2 & 1 & 1 \\
  & $(1^3)$ & 1 & 0 & 0 \\
\midrule
4 & $(4)$ & 1 & 1 & 1 \\
  & $(3,1)$ & 3 & 1 & 1 \\
  & $(2^2)$ & 2 & 1 & 1 \\
  & $(2,1^2)$ & 3 & 0 & 0 \\
  & $(1^4)$ & 1 & 0 & 0 \\
\midrule
5 & $(5)$ & 1 & 1 & 1 \\
  & $(4,1)$ & 4 & 2 & 2 \\
  & $(3,2)$ & 5 & 2 & 2 \\
  & $(3,1^2)$ & 6 & 1 & 1 \\
  & $(2^2,1)$ & 5 & 1 & 1 \\
  & $(2,1^3)$ & 4 & 0 & 0 \\
  & $(1^5)$ & 1 & 0 & 0 \\
\midrule
6 & $(6)$ & 1 & 1 & 1 \\
  & $(5,1)$ & 5 & 2 & 2 \\
  & $(4,2)$ & 9 & 3 & 3 \\
  & $(4,1^2)$ & 10 & 2 & 2 \\
  & $(3^2)$ & 5 & 2 & 2 \\
  & $(3,2,1)$ & 16 & 3 & 3 \\
  & $(3,1^3)$ & 10 & 1 & 1 \\
  & $(2^3)$ & 5 & 1 & 1 \\
  & $(2^2,1^2)$ & 9 & 0 & 0 \\
  & $(2,1^4)$ & 5 & 0 & 0 \\
  & $(1^6)$ & 1 & 0 & 0 \\
\bottomrule
\end{tabular}
\end{center}

\noindent For $n = 7$, verification confirms $r_\lambda = \dim(V_\lambda^H)$
for all 15 partitions. The totals:

\begin{center}
\begin{tabular}{c|c|c|c}
\toprule
$n$ & $\sum r_\lambda f^\lambda$ & $\sum \dim(V_\lambda^H) f^\lambda$ & $n!/2^{\lfloor n/2\rfloor}$ \\
\midrule
2 & 1 & 1 & 1 \\
3 & 3 & 3 & 3 \\
4 & 6 & 6 & 6 \\
5 & 30 & 30 & 30 \\
6 & 90 & 90 & 90 \\
7 & 630 & 630 & 630 \\
\bottomrule
\end{tabular}
\end{center}


\section{Proved corollaries}

In this section we derive the rank formula, vanishing criterion, and
hook formula \emph{assuming} Theorem~\ref{thm:main}. In subsequent
sections we prove Theorem~\ref{thm:main} for small $n$ and establish
structural lemmas toward the general case.


\subsection{Rank formula}

\begin{corollary}[Rank formula]
\label{cor:rank}
Assuming Theorem~\ref{thm:main}:
$\rk(\Pi) = n!/2^{\lfloor n/2\rfloor}$.
\end{corollary}

\begin{proof}
By the Wedderburn decomposition $\CC[\Sym_n] \cong \bigoplus_\lambda
V_\lambda^{\oplus f^\lambda}$, the total rank is
\[
  \rk(\Pi) = \sum_{\lambda \vdash n} r_\lambda \cdot f^\lambda
  = \sum_\lambda \dim(V_\lambda^H) \cdot f^\lambda.
\]
By Frobenius reciprocity,
\[
  \sum_\lambda \dim(V_\lambda^H) \cdot f^\lambda
  = \sum_\lambda \langle \mathrm{Res}^{\Sym_n}_H V_\lambda,\,
    \mathbf{1}_H \rangle \cdot \dim V_\lambda
  = \dim\!\left(\Ind_H^{\Sym_n} \mathbf{1}_H\right)
  = [\Sym_n : H] = \frac{n!}{2^{\lfloor n/2\rfloor}}. \qedhere
\]
\end{proof}

\begin{remark}
The quantity $n!/2^{\lfloor n/2\rfloor}$ is OEIS A090932, which counts
``descent-paired'' permutations: those $w \in \Sym_n$ with
$w(2i{-}1) > w(2i)$ for all $i = 1, \ldots, \lfloor n/2\rfloor$
(Callan's characterization). Computation confirms ($n \le 5$) that
$\{\Pi \cdot e_w : w \text{ descent-paired}\}$ is a basis for $\im(\Pi)$.
\end{remark}


\subsection{Vanishing criterion}

\begin{corollary}[Vanishing]
\label{cor:vanishing}
Assuming Theorem~\ref{thm:main}:
$r_\lambda = 0$ if and only if $\ell(\lambda) > \lceil(n+1)/2\rceil$.
\end{corollary}

\begin{proof}
We have $r_\lambda = 0$ iff $V_\lambda^H = \{0\}$, i.e., iff there is
no standard Young tableau $T$ of shape $\lambda$ such that the axial
distance $d_T(k) \ne \pm 1$ for all odd $k$ simultaneously in a way
that supports a $+1$ eigenvector.

More precisely, $v \in V_\lambda^H$ requires $s_k v = v$ for all odd $k$.
In Young's seminormal form, $s_k$ acts on $v_T$ with diagonal entry
$1/d$ where $d$ is the axial distance between $k$ and $k+1$ in $T$.
When $|d| = 1$, $s_k$ acts as $\pm 1$ (pure eigenvalue, no off-diagonal
mixing). When $d = +1$, $s_k v_T = v_T$, contributing to $V_\lambda^H$.
When $d = -1$, $s_k v_T = -v_T$, excluding $T$ from $V_\lambda^H$.

For a tableau $T$ to contribute to $V_\lambda^H$, we need: for each odd
$k$, entries $k$ and $k+1$ are in the same column ($d = -1$, bad) or the
same row ($d = +1$, good) or have $|d| \ge 2$ (allows mixing).

The critical constraint: if $\ell(\lambda) > \lceil(n+1)/2\rceil$,
then there are more than $\lceil(n+1)/2\rceil$ rows. The first column
has $\ell(\lambda)$ entries. Among consecutive pairs $(1,2), (3,4),
\ldots$, by pigeonhole at least one pair must occupy adjacent rows
of the first column, giving $d = -1$. This forces $s_k v_T = -v_T$
for that pair, and since there's no other tableau to mix with (swapping
would violate the SYT condition), $v_T$ is a $-1$ eigenvector of $s_k$.

Conversely, if $\ell(\lambda) \le \lceil(n+1)/2\rceil$, then entries
$1,3,5,\ldots$ can be placed in distinct rows, allowing $k$ and $k+1$
to be horizontally separated for each odd $k$. This gives
$\dim(V_\lambda^H) \ge 1$.
\end{proof}


\subsection{Hook formula}

\begin{corollary}[Hook formula]
\label{cor:hook}
Assuming Theorem~\ref{thm:main}: for hook partitions
$\lambda = (n-k, 1^k)$,
\[
  r_{(n-k,1^k)} = \binom{\lfloor(n-1)/2\rfloor}{k}.
\]
\end{corollary}

\begin{proof}
The irreducible module $V_{(n-k,1^k)} \cong \bigwedge^k V_{\mathrm{std}}$,
where $V_{\mathrm{std}}$ is the $(n{-}1)$-dimensional standard
representation (the orthogonal complement of the all-ones vector in
$\CC^n$).

\textbf{Step 1: $H$-decomposition of $V_{\mathrm{std}}$.}
Let $e_1, \ldots, e_n$ be the coordinate basis of $\CC^n$. The standard
representation has basis $\{e_i - e_{i+1} : 1 \le i \le n-1\}$ (or
any equivalent). Under $H = \langle s_1, s_3, \ldots \rangle$:
\begin{itemize}
  \item For each odd $k \le n-1$, $s_k$ swaps $e_k \leftrightarrow e_{k+1}$.
  \item The vector $e_k + e_{k+1}$ (projected to $V_{\mathrm{std}}$) is
    $H$-invariant with respect to $s_k$.
  \item The vector $e_k - e_{k+1}$ is a $-1$ eigenvector of $s_k$.
\end{itemize}

More precisely, $V_{\mathrm{std}}$ decomposes under $H$ as
$V_{\mathrm{std}} = V^H \oplus V^{-}$, where $V^H$ is the
$H$-invariant subspace and $V^{-}$ is spanned by vectors that are
$-1$ eigenvectors of at least one generator $s_k \in H$.

The $H$-invariant subspace $V_{\mathrm{std}}^H$ consists of vectors
$v = \sum a_i e_i$ (with $\sum a_i = 0$) satisfying $a_k = a_{k+1}$
for all odd $k$. This forces $a_1 = a_2$, $a_3 = a_4$, $\ldots$,
giving $\lceil n/2\rceil$ free variables minus one constraint
($\sum = 0$), so
\[
  \dim(V_{\mathrm{std}}^H) = \lceil n/2 \rceil - 1
    = \lfloor(n-1)/2\rfloor.
\]

\textbf{Step 2: $H$-invariants of the exterior power.}
We claim $(\bigwedge^k V_{\mathrm{std}})^H = \bigwedge^k(V_{\mathrm{std}}^H)$.

The inclusion $\supseteq$ is clear. For $\subseteq$: decompose
$V_{\mathrm{std}} = V^H \oplus V^{-}$ as $H$-modules. Then
\[
  \bigwedge^k V_{\mathrm{std}} = \bigoplus_{j=0}^{k}
    \bigwedge^j V^H \otimes \bigwedge^{k-j} V^{-}.
\]
An $H$-invariant element of $\bigwedge^j V^H \otimes \bigwedge^{k-j} V^-$
requires $\bigwedge^{k-j} V^-$ to contain an $H$-invariant vector.
But $V^-$ decomposes under $H$ as a direct sum of sign representations
(one for each pair generator $s_k$), so $\bigwedge^{k-j} V^-$ contains
an $H$-invariant only when $k - j = 0$, since each $s_k$ acts as $-1$
on its component of $V^-$ and hence as $(-1)^{k-j}$ on
$\bigwedge^{k-j}$ of that component.

More carefully: $V^- = \bigoplus_{i=1}^{\lfloor n/2\rfloor} L_i$ where
each $L_i$ is 1-dimensional and $s_{2i-1}$ acts as $-1$ on $L_i$ and
$+1$ on $L_j$ for $j \ne i$. Then
$\bigwedge^m V^- = \bigoplus_{|S|=m} \bigotimes_{i \in S} L_i$.
On each summand $\bigotimes_{i \in S} L_i$, the generator $s_{2i-1}$
acts as $-1$ if $i \in S$ and $+1$ otherwise. For this to be $H$-trivial,
we need $S = \emptyset$, i.e., $m = 0$.

Therefore $(\bigwedge^k V_{\mathrm{std}})^H = \bigwedge^k V^H$, and
\[
  r_{(n-k,1^k)} = \dim\!\left(\bigwedge^k V^H\right)
    = \binom{\lfloor(n-1)/2\rfloor}{k}. \qedhere
\]
\end{proof}


\section{The $\Sym_3$ lemma and structural results}

\begin{lemma}[$\Sym_3$ Lemma]
\label{lem:S3}
In any representation of $\Sym_3 = \langle s_1, s_2 \rangle$:
\[
  V^{s_1 = +1} \cap V^{s_2 = -1} = \{0\}.
\]
Equivalently, $\ker(P_2) \cap \im(P_1) = \{0\}$.
\end{lemma}

\begin{proof}
Suppose $s_1 v = v$ and $s_2 v = -v$. Then $(s_1 s_2)^3 = e$ gives
$v = (s_1 s_2)^3 v = (s_1 s_2)(s_1 s_2)(s_1 s_2) v$. But
$s_1 s_2 v = s_1(-v) = -s_1 v = -v$ (using $s_2 v = -v$ then $s_1 v = v$).
Wait --- $s_1(s_2 v) = s_1(-v) = -v$, so $(s_1 s_2)v = -v$.
Then $(s_1 s_2)^3 v = (-1)^3 v = -v$. But $(s_1 s_2)^3 = e$, so
$v = -v$, hence $v = 0$.
\end{proof}

\begin{corollary}[Within-pass injectivity]
\label{cor:injectivity}
In the staircase product, the projection $P_k$ (for $k \ge 2$) is
injective on $\im(P_{k-1})$. Hence rank drops within each block can
only occur at the $P_1$ step.
\end{corollary}

\begin{proof}
$\ker(P_k) = V^{s_k = -1}$ and $\im(P_{k-1}) \subseteq V^{s_{k-1} = +1}$.
Since $s_{k-1}$ and $s_k$ generate a copy of $\Sym_3$, the $\Sym_3$ Lemma
gives $\ker(P_k) \cap \im(P_{k-1}) = \{0\}$.
\end{proof}


\begin{corollary}
\label{cor:even-blocks}
Even-numbered blocks preserve rank. Specifically, if block $2m$ applies
$P_{2m}, P_{2m-1}, \ldots, P_1$ and the image entering block $2m$ is
contained in $V^{s_1 = +1}$, then $P_1$ acts injectively (no halving).
\end{corollary}

\begin{proof}
The image entering block $2m$ lies in $V^{s_1 = +1}$ by assumption.
The first step is $P_{2m}$, the remaining steps $P_{2m-1}, \ldots, P_2$
are injective by Corollary~\ref{cor:injectivity}. For $P_1$: the image
after $P_2$ lies in $V^{s_2 = +1} \cap V^{s_1 = +1}$ (since
$P_2$ maps $V^{s_1 = +1}$ into $V^{s_2 = +1}$ injectively, preserving
the $s_1 = +1$ condition when $s_1$ commutes with $P_2$). In fact,
the trace argument from the companion paper~\cite{rank-injectivity}
shows $P_1$ halves only when $\tr(s_1|_W) = 0$, which fails when
the image already satisfies $s_1 = +1$.
\end{proof}


\section{Block structure and the halving pattern}

\begin{proposition}[Block rank pattern]
\label{prop:block}
Computationally verified for $n \le 8$: the rank of $\Pi$ after each
complete block follows the pattern
\[
  \rk(\text{after block } k) = \frac{n!}{2^{\lceil k/2\rceil}}.
\]
That is, odd blocks halve the rank and even blocks preserve it.
After all $n-1$ blocks, $\rk(\Pi) = n!/2^{\lceil(n-1)/2\rceil}
= n!/2^{\lfloor n/2\rfloor}$.
\end{proposition}

\begin{proof}[Proof sketch]
By Corollary~\ref{cor:injectivity}, all rank drops within a block
occur at the $P_1$ step. Block $k$ applies $P_k, P_{k-1}, \ldots, P_1$.
The first projection $P_k$ maps into $V^{s_k = +1}$; the cascade
$P_{k-1}, \ldots, P_2$ is injective (shifting the eigenspace condition
from $s_k$ to $s_{k-1}$ to $\ldots$ to $s_2$); then $P_1$ either
halves (if the image has zero $s_1$-trace) or is injective (if the
image already lies in $V^{s_1 = +1}$).

The halving at odd blocks and preservation at even blocks is proved
for blocks 1--2 unconditionally (see the companion paper), and
the trace-preservation mechanism extends to later blocks under
verified conditions.
\end{proof}


\section{Proof of the main theorem: small cases}

\subsection{$n = 2$}

$\Sym_2$ has two irreps: $V_{(2)}$ (trivial, $f = 1$) and $V_{(1^2)}$
(sign, $f = 1$). The staircase word is $(s_1)$, so $\Pi = P_1$.
\begin{itemize}
  \item $V_{(2)}$: $s_1 = +1$, so $P_1 = 1$, $r_{(2)} = 1$.
    $H = \langle s_1 \rangle$, $V_{(2)}^H = V_{(2)}$, $\dim = 1$. \checkmark
  \item $V_{(1^2)}$: $s_1 = -1$, so $P_1 = 0$, $r_{(1^2)} = 0$.
    $V_{(1^2)}^H = \{0\}$, $\dim = 0$. \checkmark
\end{itemize}


\subsection{$n = 3$}

$\Sym_3$ has three irreps: $(3), (2,1), (1^3)$ with $f = 1, 2, 1$.
$H = \langle s_1 \rangle$, staircase word $(s_1, s_2, s_1)$.

\begin{itemize}
  \item $V_{(3)}$: trivial. $\Pi = 1$. $r = 1 = \dim V^H$. \checkmark
  \item $V_{(1^3)}$: sign. $s_1 = s_2 = -1$. $P_1 = 0$, so $\Pi = 0$.
    $r = 0$. $V^H = V^{s_1=+1} = \{0\}$ since $s_1 = -1$. \checkmark
  \item $V_{(2,1)}$: standard 2-dimensional representation. In the
    seminormal basis, $s_1 = \begin{pmatrix} 1 & 0 \\ 0 & -1\end{pmatrix}$,
    $s_2 = \begin{pmatrix} -1/2 & \sqrt{3}/2 \\ \sqrt{3}/2 & 1/2\end{pmatrix}$.
    Then $P_1 = \begin{pmatrix} 1 & 0 \\ 0 & 0\end{pmatrix}$,
    $P_2 = \begin{pmatrix} 1/4 & \sqrt{3}/4 \\ \sqrt{3}/4 & 3/4\end{pmatrix}$.
    Computing $\Pi = P_1 P_2 P_1$: $P_2 P_1 = \begin{pmatrix} 1/4 & 0 \\
    \sqrt{3}/4 & 0\end{pmatrix}$, then $P_1(P_2 P_1) = \begin{pmatrix}
    1/4 & 0 \\ 0 & 0\end{pmatrix}$.
    So $r_{(2,1)} = 1$. Meanwhile $V^H = V^{s_1=+1} = \mathrm{span}(e_1)$,
    $\dim = 1$. \checkmark
\end{itemize}


\section{Toward the general proof}

\subsection{Upper bound: $r_\lambda \le \dim(V_\lambda^H)$}

\begin{proposition}
\label{prop:upper}
For all $\lambda \vdash n$: $\im(\Pi_\lambda) \subseteq V_\lambda^{s_1}$
(the $+1$-eigenspace of $s_1$). Hence $r_\lambda \le \dim(V_\lambda^{s_1})$.
\end{proposition}

\begin{proof}
The rightmost factor in the staircase product is $P_1$, so
$\im(\Pi_\lambda) \subseteq \im(P_1^\lambda) = V_\lambda^{s_1}$.
\end{proof}

\begin{remark}
This gives $r_\lambda \le \dim(V_\lambda^{s_1})$, which is weaker
than $r_\lambda \le \dim(V_\lambda^H)$ since $V_\lambda^H \subseteq
V_\lambda^{s_1}$. However, computationally, equality
$r_\lambda = \dim(V_\lambda^H)$ holds, and the stronger upper bound
$\im(\Pi_\lambda) \subseteq V_\lambda^H$ does \emph{not} hold
(the image is not contained in $V_\lambda^H$ in general --- the
product $\Pi$ does not absorb the pair transpositions $s_3, s_5, \ldots$).
\end{remark}


\subsection{The equality: what remains}

The main theorem asserts rank equality $r_\lambda = \dim(V_\lambda^H)$
without asserting containment of the image in $V_\lambda^H$. The
computational evidence shows that $\im(\Pi_\lambda) \ne V_\lambda^H$
in general (for $n \ge 4$, $\Pi$ does not absorb $s_3, s_5, \ldots$),
yet the \emph{dimensions} match.

The proof strategy from the companion paper~\cite{rank-injectivity}
establishes this via the halving mechanism:
\begin{enumerate}
  \item Within-pass injectivity: only $P_1$ steps can drop rank
    (Corollary~\ref{cor:injectivity}).
  \item Odd blocks halve via vanishing $s_1$-trace, even blocks preserve
    (Proposition~\ref{prop:block}).
  \item After $\lfloor n/2\rfloor$ halvings, the rank is
    $n!/2^{\lfloor n/2\rfloor}$ in the regular representation.
\end{enumerate}

The per-irrep statement $r_\lambda = \dim(V_\lambda^H)$ requires
tracking the halving in each irreducible component, which is the
remaining gap. The halving mechanism must show that in each $V_\lambda$,
exactly $\dim(V_\lambda^H)$ dimensions survive all $\lfloor n/2\rfloor$
halvings.


\subsection{Connection to the Solomon descent algebra}

The staircase word interacts with the descent structure of $\Sym_n$.
The Solomon descent algebra $\Sigma_n \subseteq \CC[\Sym_n]$ has
basis elements $D_S = \sum_{w : \mathrm{Des}(w) = S} w$ indexed by
subsets $S \subseteq \{1, \ldots, n-1\}$. The product $\Pi$ is an
element of $\CC[\Sym_n]$ that respects the descent filtration in a
way that produces the A090932 counts. A deeper understanding of this
connection --- particularly the fact that descent-paired permutations
(a descent-type condition) index the image basis --- may yield a
proof of the per-irrep statement.


\section{Summary}

\begin{center}
\begin{tabular}{l|l}
\toprule
\textbf{Result} & \textbf{Status} \\
\midrule
$r_\lambda = \dim(V_\lambda^H)$ & Verified $n \le 7$; proved $n \le 3$ \\
$\rk(\Pi) = n!/2^{\lfloor n/2\rfloor}$ & Corollary (assuming main thm) \\
Vanishing: $r_\lambda = 0 \Leftrightarrow \ell(\lambda) > \lceil(n+1)/2\rceil$
  & Corollary (assuming main thm) \\
Hook: $r_{(n-k,1^k)} = \binom{\lfloor(n-1)/2\rfloor}{k}$
  & Corollary (assuming main thm) \\
$\Sym_3$ Lemma & Proved \\
Within-pass injectivity & Proved \\
Block halving pattern & Verified $n \le 8$ \\
First two halvings & Proved (companion paper) \\
\bottomrule
\end{tabular}
\end{center}

The main open problem is to prove $r_\lambda = \dim(V_\lambda^H)$ for
general $\lambda$ and $n$. The halving mechanism (companion paper)
establishes the total rank but the per-irrep refinement remains.


\begin{thebibliography}{9}
\bibitem{rank-injectivity}
Clio, \emph{The Rank of the Symmetrizing Product: Injectivity, Cascade
Surjectivity, and the Formula $n!/2^{\lfloor n/2\rfloor}$}, April 14, 2026.

\bibitem{Galashin2020}
P.~Galashin, \emph{Symmetries of stochastic colored vertex models},
Ann.\ Probab.\ \textbf{49} (2021), no.~5, 2175--2219.

\bibitem{OEIS-A090932}
OEIS Foundation, \emph{A090932: $n!/2^{\lfloor n/2\rfloor}$},
\texttt{https://oeis.org/A090932}.
\end{thebibliography}

\end{document}
